% =====================================================================
%  Chuong99_KetLuan.tex — Kết luận, Tài liệu tham khảo, Phụ lục
%
%  Tách từ main.tex (dòng 2548-2601 trong bản cũ).
%  Chứa:
%    - Phần Kết luận (Tóm tắt, Hạn chế, Hướng phát triển)
%    - Tài liệu tham khảo (BibTeX)
%    - Phụ lục (qua % =====================================================================
%  PHỤ LỤC -- các chi tiết bổ sung không gói gọn trong các chương chính
% =====================================================================

\section*{Phụ lục A: Thống kê chi tiết bộ dữ liệu}

\subsection*{A.1. Số mẫu theo lớp của từng tập}

\begin{table}[H]
  \centering
  \small
  \caption[Số mẫu theo lớp của từng tập]{Số mẫu theo lớp (Thật / Giả) trong ba tập sau khi chia phân tầng với \texttt{random\_state=42}.}
  \label{tab:appendix_split_class}
  \begin{tabular}{l r r r}
    \toprule
    \textbf{Tập} & \textbf{\textit{thật} (0)} & \textbf{\textit{giả} (1)} & \textbf{Tổng} \\
    \midrule
    Train & \ntrainreal{} & \ntrainreal{} & \ntrain{} \\
    Val   & \nvalreal & \nvalfake & \nval{} \\
    Test  & \ntestreal & \ntestfake & \ntest{} \\
    \bottomrule
  \end{tabular}
\end{table}

Sau khi chia tập, nguy cơ rò rỉ dữ liệu được kiểm tra theo cả trường \texttt{text} và trường \texttt{id}; kết quả cho thấy cả hai phép giao đều có kích thước bằng~0, đảm bảo tính độc lập giữa các tập.

\subsection*{A.2. Ví dụ mẫu trong tập kiểm tra}

\noindent\textbf{Mẫu tin giả (label = 1, ngắn):}
\begin{quote}\small\itshape
Học\_sinh có\_thể nhớ bài tốt hơn nếu đặt sách\_giáo\_khoa dưới
gối trong lúc ngủ.\hfill(id=9274)
\end{quote}

\noindent\textbf{Mẫu tin giả (label = 1, dài):}
\begin{quote}\small\itshape
``Một thử\_nghiệm tại trường\_học cho thấy việc nghe nhạc bằng
tai\_nghe trong 15 phút trước khi kiểm\_tra có\_thể làm tăng khả\_năng
ghi\_nhớ của học\_sinh lên gần 40\%. Các giáo\_viên vì vậy được
khuyến\_nghị cho học\_sinh nghe nhạc trước mỗi bài\_thi\ldots''
\hfill(id=6385, Viện Nghiên cứu Giáo dục, về phương pháp học tập)
\end{quote}


\noindent\textbf{Mẫu tin giả (label = 1):}
\begin{quote}\small\itshape
Tầng ozone đã hoàn\_toàn biến mất vào năm 2023 do sự nóng lên
của mặt\_trời.\hfill(id=8262)
\end{quote}

\noindent\textbf{Mẫu tin giả giật tít (label = 1):}
\begin{quote}\small\itshape
Chính\_quyền cấp xã luôn đi\_tắt\_đón\_đầu dân. Dân chờ cả ngày
trời để lãnh ``hỗ\_trợ'' 80 ngàn đồng rồi bị trấn\_lột tại\_chỗ.\hfill(id=3126)
\end{quote}

\subsection*{A.3. Pipeline tiền xử lý}

Quy trình làm sạch bao gồm bốn bước theo thứ tự: (i) loại bỏ các bản ghi trùng lặp dựa trên cột \texttt{text}; (ii) loại bỏ những bản ghi chứa ít hơn \texttt{min\_word\_count = \minwordcount{}} từ; (iii) chuẩn hóa cột \texttt{date} về định dạng ISO~8601; (iv) đánh lại chỉ mục \texttt{id} tuần tự từ~1.

Phân đoạn từ ưu tiên sử dụng \texttt{py\_vncorenlp} (RDRSegmenter --- cùng tokenizer với PhoBERT); hệ thống tự động chuyển sang \texttt{underthesea.word\_tokenize} trong trường hợp thư viện chính không khả dụng.

\section*{Phụ lục B: Cấu hình huấn luyện chi tiết}

Toàn bộ siêu tham số được tập trung tại file cấu hình chung của dự án, được đọc thông qua singleton \texttt{cfg}.

\begin{table}[H]
  \centering
  \small
  \caption[Cấu hình TF-IDF, LR, SVM]{Cấu hình TF-IDF, Logistic Regression và SVM.}
  \label{tab:appendix_tfidf_lr_svm}
  \begin{tabular}{l l l}
    \toprule
    \textbf{Nhóm} & \textbf{Siêu tham số} & \textbf{Giá trị} \\
    \midrule
    \multirow{4}{*}{TF-IDF}
      & \texttt{max\_features}     & \tfidfmaxfeatures{} \\
      & \texttt{ngram\_range}      & $(1, 2)$ \\
      & \texttt{min\_df / max\_df} & 5 / 0{,}95 \\
      & \texttt{sublinear\_tf}     & True \\
    \midrule
    \multirow{4}{*}{LR}
      & \texttt{C} (tốt nhất)      & \lrbestc{} (solver=saga, L2) \\
      & \texttt{max\_iter}         & 3000 \\
      & \texttt{class\_weight}     & balanced \\
      & Lưới tìm kiếm             & $C \in \{0{,}5;\ 1;\ 5;\ 10;\ 20;\ 50\}$ \\
    \midrule
    \multirow{3}{*}{SVM}
      & \texttt{C} (tốt nhất)      & \svmbestc{} (LinearSVC + Platt) \\
      & \texttt{class\_weight}     & balanced \\
      & Lưới tìm kiếm             & $C \in \{0{,}1;\ 0{,}5;\ 1;\ 2;\ 5;\ 10\}$ \\
    \bottomrule
  \end{tabular}
\end{table}

\begin{table}[H]
  \centering
  \small
  \caption[Cấu hình BiLSTM và PhoBERT]{Cấu hình BiLSTM và PhoBERT.}
  \label{tab:appendix_dl}
  \begin{tabular}{l l l}
    \toprule
    \textbf{Mô hình} & \textbf{Siêu tham số} & \textbf{Giá trị} \\
    \midrule
    \multirow{7}{*}{BiLSTM}
      & Embedding dim            & \bilstmembdim{} (FastText cc.vi.300) \\
      & Hidden dim / Layers      & \bilstmhiddendim{} / \bilstmlayers{} \\
      & Max seq length           & \bilstmmaxseq{} \\
      & Dropout / Weight decay   & 0{,}3 / $10^{-4}$ \\
      & Learning rate            & $10^{-3}$ (AdamW) \\
      & Batch / Epochs           & 64 / 20 (patience=5) \\
      & Gradient clip            & $\max\_norm = 1{,}0$ \\
    \midrule
    \multirow{7}{*}{PhoBERT}
      & \texttt{model\_name}     & \texttt{vinai/phobert-base} \\
      & Max seq length           & \phobertmaxseq{} \\
      & Dropout                  & 0{,}1 \\
      & LR / Weight decay        & $3 \times 10^{-5}$ / 0{,}01 \\
      & Batch / Grad accum.      & 16 / 4 (effective $\approx 64$) \\
      & Epochs / Patience        & 8 / 2 \\
      & Layer LR decay           & 0{,}95 \\
    \bottomrule
  \end{tabular}
\end{table}

\subsection*{B.1. Knowledge Distillation}

Mô hình sinh viên là \texttt{StudentBiLSTM} --- một phiên bản BiLSTM rút gọn với \texttt{embedding\_dim = \studentembdim{}}, \texttt{hidden\_dim = \studenthiddim{}}, \texttt{num\_layers = \studentlayers{}}. Trọng số soft target $\alpha = \studentalpha{}$, nhiệt độ phân phối $T = \studenttemp{}$. Optimizer sử dụng AdamW với \texttt{lr}=$10^{-3}$, scheduler CosineAnnealingLR. Huấn luyện tối đa \studentepochs{} epoch với early stopping (\texttt{patience}=\studentpatience{}). Hàm mất mát KD chuẩn (Hinton et al., 2015) đã được trình bày ở Mục~\ref{sec:extended_analyses}; mô hình sinh viên được huấn luyện trên tập logits giáo viên đã được tính trước và lưu lại nhằm tiết kiệm chi phí tính toán.

\subsection*{B.2. Post-hoc Calibration}

Ba phương pháp hiệu chỉnh xác suất --- đều được khớp trên tập validation và áp dụng lên tập kiểm tra:

\begin{itemize}[leftmargin=2em,itemsep=2pt]
  \item \textbf{Platt scaling:} hồi quy logistic một chiều trên logits.
  \item \textbf{Temperature scaling:} tối ưu vô hướng $T \in [10^{-3},\ 10^{3}]$ bằng thuật toán L-BFGS-B.
  \item \textbf{Isotonic regression:} ánh xạ đơn điệu phi tham số trên $\sigma(z)$.
\end{itemize}

Giá trị $T$ đã khớp tương ứng: LR $=\lrtfit{}$, SVM $=\svmtfit{}$, BiLSTM $=\bilstmtfit{}$, PhoBERT $=\phoberttfit{}$.

\section*{Phụ lục C: Kết quả chi tiết}

\subsection*{C.1. Metric tổng hợp trên tập test}

\begin{table}[H]
  \centering
  \small
  \caption[Metric tổng hợp trên tập test]{Accuracy, Precision, Recall và F1-macro trên tập test.}
  \label{tab:appendix_full_metrics}
  \begin{tabular}{l c c c c}
    \toprule
    \textbf{Mô hình} & \textbf{Accuracy} & \textbf{Precision} & \textbf{Recall} & \textbf{F1-macro} \\
    \midrule
    Logistic Regression & \lracctab{} & \lrprectab{} & \lrrectab{} & \lrfonetab{} \\
    SVM                 & \svmacctab{} & \svmprectab{} & \svmrectab{} & \svmfonetab{} \\
    BiLSTM              & \bilstmacctab{} & \bilstmprectab{} & \bilstmrectab{} & \bilstmfonetab{} \\
    PhoBERT             & \phobertacctab{} & \phobertprectab{} & \phobertrectab{} & \phobertfonetab{} \\
    \bottomrule
  \end{tabular}
\end{table}

\subsection*{C.2. Metric theo lớp}

\begin{table}[H]
  \centering
  \small
  \caption[Precision/Recall/F1 theo lớp]{Precision, Recall và F1 theo lớp (Thật / Giả) trên tập test.}
  \label{tab:appendix_per_class}
  \begin{tabular}{l c c c c c c}
    \toprule
    & \multicolumn{3}{c}{\textbf{Lớp Thật (0)}} & \multicolumn{3}{c}{\textbf{Lớp Giả (1)}} \\
    \cmidrule(lr){2-4} \cmidrule(lr){5-7}
    \textbf{Mô hình} & Prec & Rec & F1 & Prec & Rec & F1 \\
    \midrule
    Logistic Regression & \lrrealprec{} & \lrrealrec{} & \lrrealfone{} & \lrfakeprec{} & \lrfakerec{} & \lrfakefone{} \\
    SVM                 & \svmrealprec{} & \svmrealrec{} & \svmrealfone{} & \svmfakeprec{} & \svmfakerec{} & \svmfakefone{} \\
    BiLSTM              & \bilstmrealprec{} & \bilstmrealrec{} & \bilstmrealfone{} & \bilstmfakeprec{} & \bilstmfakerec{} & \bilstmfakefone{} \\
    PhoBERT             & \phobertrealprec{} & \phobertrealrec{} & \phobertrealfone{} & \phobertfakeprec{} & \phobertfakerec{} & \phobertfakefone{} \\
    \bottomrule
  \end{tabular}
\end{table}

\subsection*{C.3. Confusion matrix trên tập test}

\begin{table}[H]
  \centering
  \small
  \caption[Confusion matrix trên tập test]{Số mẫu TN, FP, FN, TP của bốn mô hình trên tập test.}
  \label{tab:appendix_confusion}
  \begin{tabular}{l r r r r r}
    \toprule
    \textbf{Mô hình} & \textbf{TN} & \textbf{FP} & \textbf{FN} & \textbf{TP} & \textbf{Tổng lỗi} \\
    \midrule
    Logistic Regression & 985  & \lrfpcount{}   & \lrfncount{}   & 992 & \lrerrors{} (\lrerrorrate\%) \\
    SVM                 & 1.010 & \svmfpcount{}  & \svmfncount{}  & 756 & \svmerrors{} (\svmerrorrate\%) \\
    BiLSTM              & 977  & \bilstmfpcount{} & \bilstmfncount{} & 751 & \bilstmerrors{} (\bilstmerrorrate\%) \\
    PhoBERT             & \phoberttN{} & \phobertfP{} & \phobertfN{} & \phoberttP{} & \phoberterrors{} (\phoberterrorrate\%) \\
    \bottomrule
  \end{tabular}
\end{table}

\subsection*{C.4. Khoảng tin cậy Bootstrap 95\%}

Khoảng tin cậy được đo bằng phương pháp bootstrap với 10{,}000 lần lặp.

\begin{table}[H]
  \centering
  \small
  \caption[Khoảng tin cậy Bootstrap 95\% cho F1-score]{Khoảng tin cậy 95\% cho F1-score trên tập test, ước lượng bằng bootstrap với 10{,}000 lần lặp.}
  \label{tab:appendix_bootstrap}
  \begin{tabular}{l c}
    \toprule
    \textbf{Mô hình} & \textbf{95\% CI (F1-score)} \\
    \midrule
    Logistic Regression & \lrfonecI{} \\
    SVM                 & \svmfonecI{} \\
    BiLSTM              & \bilstmfonecI{} \\
    PhoBERT             & \phobertfonecI{} \\
    \bottomrule
  \end{tabular}
\end{table}

\subsection*{C.5. Knowledge Distillation -- chi tiết}

\begin{table}[H]
  \centering
  \small
  \caption[So sánh giáo viên -- sinh viên]{So sánh hai mô hình giáo viên (PhoBERT, BiLSTM) và sinh viên BiLSTM ($\alpha=\studentalpha$, $T=\studenttemp$).}
  \label{tab:appendix_distillation}
  \begin{tabular}{l r r r}
    \toprule
    \textbf{Mô hình} & \textbf{F1 test} & \textbf{Params} & \textbf{Kích thước (MB)} \\
    \midrule
    PhoBERT (giáo viên)  & \phobertfonetab{} & \phobertparamcount{} & \phobertsizemB{} \\
    BiLSTM (giáo viên)   & \bilstmfonetab{}  & \bilstmparamcount{}  & 24{,}57 \\
    Student BiLSTM (KD)  & \studentfone{}    & \studentparams{}    & \studentsizemB{} \\
    \bottomrule
  \end{tabular}
\end{table}

\noindent Tỷ lệ nén so với PhoBERT là \studentcompressionphi{$\times$}; so với BiLSTM là \studentcompressionbI{$\times$}. Tốc độ suy diễn đạt \studentlat{}~ms, nhanh hơn \studentspeedup{$\times$}; mô hình sinh viên giữ được \studentretainbI{} phần trăm chất lượng của BiLSTM.

\subsection*{C.6. Attribution Faithfulness}

Faithfulness là mức giảm xác suất trung bình khi che \texttt{top-5} token được attribution đánh dấu trên tập \nattrexamples{} mẫu.

\begin{table}[H]
  \centering
  \small
  \caption[Faithfulness của sáu phương pháp attribution]{Faithfulness của sáu phương pháp attribution.}
  \label{tab:appendix_faithfulness}
  \begin{tabular}{l c}
    \toprule
    \textbf{Phương pháp} & \textbf{Faithfulness} \\
    \midrule
    LR -- SHAP                     & \attrfaithfulnesslR{} \\
    SVM -- SHAP                    & \attrfaithfulnesssVM{} \\
    BiLSTM -- Integrated Gradients & \attrfaithfulnessbilstm{} \\
    BiLSTM -- Simple Gradients     & \attrfaithfulnessbilstmgrad{} \\
    PhoBERT -- SHAP                & \attrfaithfulnessphobERTSHAP{} \\
    PhoBERT -- Integrated Gradients& \attrfaithfulnessphobERTIG{} \\
    PhoBERT -- Attention Rollout   & \attrfaithfulnessphobERTRollout{} \\
    \bottomrule
  \end{tabular}
\end{table}

\subsection*{C.7. Cross-validation 5-fold}

\begin{table}[H]
  \centering
  \small
  \caption[Cross-validation 5-fold]{Accuracy và F1 trung bình qua 5-fold cross-validation (3 seed) cho hai mô hình truyền thống.}
  \label{tab:appendix_cv}
  \begin{tabular}{l c c c c}
    \toprule
    \textbf{Mô hình} & \textbf{Accuracy} & \textbf{$\pm$std} & \textbf{F1-macro} & \textbf{$\pm$std} \\
    \midrule
    Logistic Regression & \cvlRAcc{} & \cvlRAccstd{} & \cvlRFone{} & \cvlRFonestd{} \\
    SVM                 & \cvsVMAcc{} & \cvsVMAccstd{} & \cvsVMFone{} & \cvsVMFonestd{} \\
    \bottomrule
  \end{tabular}
\end{table}

\subsection*{C.8. Hard cases -- phân bố}

Tập \nhardexamples{} mẫu khó (cả bốn mô hình đều phân loại sai) được phân bố theo lý do thất bại suy đoán như sau:

\begin{table}[H]
  \centering
  \small
  \caption[Phân bố hard cases theo lý do]{Phân bố \nhardexamples{} mẫu khó theo lý do thất bại suy đoán.}
  \label{tab:appendix_hard_reasons}
  \begin{tabular}{l r}
    \toprule
    \textbf{Lý do} & \textbf{Số mẫu} \\
    \midrule
    \texttt{dataset\_ambiguity}      & \nhardambiguity{} \\
    \texttt{knowledge\_verification} & \nhardknowledge{} \\
    \texttt{semantic\_reasoning}     & \nhardsemantic{} \\
    \midrule
    \textbf{Chủ đề politics}        & \nhardpolitics{} (phân bố theo chủ đề) \\
    \textbf{Độ dài short ($<$30 từ)} & \nhardshort{} \\
    \bottomrule
  \end{tabular}
\end{table}

\section*{Phụ lục D: Tái lập thực nghiệm}

\subsection*{D.1. Yêu cầu môi trường}

\noindent\begin{tabular}{@{}l@{\;}l@{}}
\textbf{Python}     & \texttt{numpy 2.4}, \texttt{pandas 3.0}, \texttt{scikit-learn 1.8}, \\
                    & \texttt{scipy 1.17}, \texttt{torch 2.10}, \texttt{transformers 4.57}, \\
                    & \texttt{sentencepiece 0.2}, \texttt{fasttext-wheel 0.9.2}. \\
\textbf{Tiếng Việt} & \texttt{py\_vncorenlp 0.1.4} (Java $\geq$ 11) \\
                    & hoặc \texttt{underthesea 9.1} (dự phòng). \\
\textbf{Giải thích} & \texttt{shap $\geq$ 0.45}, \texttt{captum $\geq$ 0.7}. \\
\end{tabular}

Biến môi trường cần thiết: \texttt{JAVA\_HOME} trỏ tới JRE. PhoBERT (\texttt{vinai/phobert-base}) và FastText (\texttt{cc.vi.300.bin}) được tải tự động từ HuggingFace và lưu tại thư mục \texttt{data/fasttext/}.

\subsection*{D.2. Pipeline huấn luyện -- đánh giá}

\begin{verbatim}
python -m venv .venv && source .venv/bin/activate
pip install -r requirements.txt
python -m src.preprocessing.word_segmentation
python -m src.preprocessing.split_data
python -m src.features.extract_all_features
python -m src.training.train_all
python -m src.training.reproduce_predictions
python -m src.training.train_student
python -m src.evaluation.post_hoc_calibration
python -m src.evaluation.error_analysis
python -m src.evaluation.hard_cases
python -m src.analysis.statistical_tests
python -m src.analysis.generate_paper_tables
python -m src.analysis.generate_paper_figures
cd paper && ./build.sh
\end{verbatim}

\subsection*{D.3. Thiết lập ngẫu nhiên}

Seed chính của toàn bộ pipeline là \texttt{cfg.RANDOM\_STATE = \randomstate{}} --- được sử dụng cho \texttt{train\_test\_split}, \texttt{GridSearchCV} (qua \texttt{StratifiedKFold(shuffle=True)}) và các script huấn luyện. Hàm \texttt{set\_reproducibility\_seeds} cố định đồng thời các seed của \texttt{random}, \texttt{numpy}, \texttt{torch} và \texttt{torch.cuda}; hàm này được gọi ngay đầu mỗi trainer nhằm đảm bảo kết quả hoàn toàn tái lập.

\subsection*{D.4. File cấu hình}

Toàn bộ siêu tham số được tập trung tại file cấu hình chung và import qua singleton \texttt{cfg}, bao gồm các khối: \texttt{PATHS}, \texttt{DATA} (tỉ lệ chia tập, \texttt{min\_word\_count = \minwordcount{}}), \texttt{TFIDF}, \texttt{LR}, \texttt{SVM}, \texttt{BILSTM}, \texttt{PHOBERT}, \texttt{ANALYSIS} (\texttt{bootstrap\_iterations = 10{,}000}, \texttt{significance\_level = 0{,}05}).

Mọi giá trị hiển thị trong bài báo được định nghĩa tại \texttt{data\_commands.sty} dưới dạng \texttt{\textbackslash newcommand}, cho phép cập nhật đồng bộ từ một vị trí duy nhất --- tránh sai lệch giữa các phần của tài liệu.
)
%
%  Lưu ý: Phần này dùng \chapter*{} (không đánh số) vì đã có Chương 5 (Kết
%  quả thực nghiệm) là chương nội dung cuối cùng. Kết luận và Phụ lục là
%  các phần bổ sung theo template đồ án.
% =====================================================================

\chapter*{Kết luận}
\addcontentsline{toc}{chapter}{Kết luận}

\section*{Tóm tắt đóng góp}

Nghiên cứu đã xây dựng và đối sánh bốn phương pháp phát hiện tin giả tiếng Việt trên bộ dữ liệu \nclean{} mẫu sau khi làm sạch. Những kết quả đáng chú ý nhất bao gồm:

\begin{itemize}[leftmargin=2em]
  \item PhoBERT đạt hiệu năng vượt trội với accuracy = \phobertacc\%, F1 = \phobertfone\%, AUC = \phobertaUC.
  \item Mức cải thiện so với baseline Logistic Regression là +\deltalR{} điểm F1.
  \item Kiểm định McNemar với hiệu chỉnh Holm--Bonferroni xác nhận sự vượt trội có ý nghĩa thống kê ($p < 0{,}001$).
  \item SVM vượt BiLSTM --- minh chứng rõ ràng cho thấy chất lượng đặc trưng đóng vai trò quan trọng hơn độ phức tạp của kiến trúc mô hình.
  \item Nghiên cứu ablation cho thấy tiền huấn luyện trên ngữ liệu tiếng Việt là nhân tố quyết định; phân đoạn từ ảnh hưởng không đáng kể đối với TF-IDF nhưng lại rất quan trọng đối với PhoBERT.
  \item Phân tích lỗi chỉ ra \nhardexamples{} mẫu khó --- những mẫu mà cả bốn mô hình đều phân loại sai --- cho thấy hệ thống cần được bổ sung các nguồn tri thức bên ngoài mới có thể xử lý hiệu quả hơn.
\end{itemize}

\section*{Hạn chế}

\begin{itemize}[leftmargin=2em]
  \item Tập dữ liệu hiện có thể chưa phản ánh đầy đủ mọi thể loại và chiến thuật lan truyền tin giả đang hiện hữu trong thực tế, đặc biệt với các dạng nội dung mới xuất hiện gần đây hoặc những nội dung mang tính chuyên biệt theo từng lĩnh vực.

  \item Hiệu năng mô hình có thể bị suy giảm đáng kể khi áp dụng lên dữ liệu nằm ngoài miền đã huấn luyện (domain shift), nhất là đối với những hình thức ngụy trang tinh vi hoặc các biến thể ngôn ngữ chưa từng gặp.

  \item PhoBERT yêu cầu tài nguyên tính toán tương đối lớn --- đây là rào cản đáng kể khi triển khai ở quy mô lớn hoặc trong các bài toán đòi hỏi suy luận thời gian thực.
\end{itemize}

\section*{Hướng phát triển}

\begin{itemize}[leftmargin=2em]

  \item Tiếp tục mở rộng nguồn dữ liệu bằng cách thu thập các bài viết tiếng Việt có nội dung và lĩnh vực đa dạng, đồng thời cập nhật những dạng tin giả mới xuất hiện trong thực tế.

  \item Bổ sung các đặc trưng ngoài văn bản, chẳng hạn như hình ảnh và thông tin về quá trình lan truyền trên mạng xã hội, nhằm hỗ trợ mô hình trong quá trình nhận diện tin giả.

  \item Hướng đến khả năng phát hiện tin giả theo thời gian thực bằng cách triển khai một phiên bản PhoBERT có kích thước gọn hơn nhờ áp dụng kỹ thuật knowledge distillation.

  \item Đưa quy trình fact-checking vào kết hợp với kết quả phân loại của mô hình, qua đó nâng cao mức độ đáng tin cậy của hệ thống khi đưa ra kết luận.

\end{itemize}
% ============================================================
%  Tài liệu tham khảo (BibTeX)
% ============================================================
\bibliographystyle{ieeetr}
\bibliography{refs}

% ============================================================
%  Phụ lục (tách sang file front/pages/PhuLuc.tex)
% ============================================================
\cleardoublepage
\chapter*{Phụ lục}
\addcontentsline{toc}{chapter}{Phụ lục}
% =====================================================================
%  PHỤ LỤC -- các chi tiết bổ sung không gói gọn trong các chương chính
% =====================================================================

\section*{Phụ lục A: Thống kê chi tiết bộ dữ liệu}

\subsection*{A.1. Số mẫu theo lớp của từng tập}

\begin{table}[H]
  \centering
  \small
  \caption[Số mẫu theo lớp của từng tập]{Số mẫu theo lớp (Thật / Giả) trong ba tập sau khi chia phân tầng với \texttt{random\_state=42}.}
  \label{tab:appendix_split_class}
  \begin{tabular}{l r r r}
    \toprule
    \textbf{Tập} & \textbf{\textit{thật} (0)} & \textbf{\textit{giả} (1)} & \textbf{Tổng} \\
    \midrule
    Train & \ntrainreal{} & \ntrainreal{} & \ntrain{} \\
    Val   & \nvalreal & \nvalfake & \nval{} \\
    Test  & \ntestreal & \ntestfake & \ntest{} \\
    \bottomrule
  \end{tabular}
\end{table}

Sau khi chia tập, nguy cơ rò rỉ dữ liệu được kiểm tra theo cả trường \texttt{text} và trường \texttt{id}; kết quả cho thấy cả hai phép giao đều có kích thước bằng~0, đảm bảo tính độc lập giữa các tập.

\subsection*{A.2. Ví dụ mẫu trong tập kiểm tra}

\noindent\textbf{Mẫu tin giả (label = 1, ngắn):}
\begin{quote}\small\itshape
Học\_sinh có\_thể nhớ bài tốt hơn nếu đặt sách\_giáo\_khoa dưới
gối trong lúc ngủ.\hfill(id=9274)
\end{quote}

\noindent\textbf{Mẫu tin giả (label = 1, dài):}
\begin{quote}\small\itshape
``Một thử\_nghiệm tại trường\_học cho thấy việc nghe nhạc bằng
tai\_nghe trong 15 phút trước khi kiểm\_tra có\_thể làm tăng khả\_năng
ghi\_nhớ của học\_sinh lên gần 40\%. Các giáo\_viên vì vậy được
khuyến\_nghị cho học\_sinh nghe nhạc trước mỗi bài\_thi\ldots''
\hfill(id=6385, Viện Nghiên cứu Giáo dục, về phương pháp học tập)
\end{quote}


\noindent\textbf{Mẫu tin giả (label = 1):}
\begin{quote}\small\itshape
Tầng ozone đã hoàn\_toàn biến mất vào năm 2023 do sự nóng lên
của mặt\_trời.\hfill(id=8262)
\end{quote}

\noindent\textbf{Mẫu tin giả giật tít (label = 1):}
\begin{quote}\small\itshape
Chính\_quyền cấp xã luôn đi\_tắt\_đón\_đầu dân. Dân chờ cả ngày
trời để lãnh ``hỗ\_trợ'' 80 ngàn đồng rồi bị trấn\_lột tại\_chỗ.\hfill(id=3126)
\end{quote}

\subsection*{A.3. Pipeline tiền xử lý}

Quy trình làm sạch bao gồm bốn bước theo thứ tự: (i) loại bỏ các bản ghi trùng lặp dựa trên cột \texttt{text}; (ii) loại bỏ những bản ghi chứa ít hơn \texttt{min\_word\_count = \minwordcount{}} từ; (iii) chuẩn hóa cột \texttt{date} về định dạng ISO~8601; (iv) đánh lại chỉ mục \texttt{id} tuần tự từ~1.

Phân đoạn từ ưu tiên sử dụng \texttt{py\_vncorenlp} (RDRSegmenter --- cùng tokenizer với PhoBERT); hệ thống tự động chuyển sang \texttt{underthesea.word\_tokenize} trong trường hợp thư viện chính không khả dụng.

\section*{Phụ lục B: Cấu hình huấn luyện chi tiết}

Toàn bộ siêu tham số được tập trung tại file cấu hình chung của dự án, được đọc thông qua singleton \texttt{cfg}.

\begin{table}[H]
  \centering
  \small
  \caption[Cấu hình TF-IDF, LR, SVM]{Cấu hình TF-IDF, Logistic Regression và SVM.}
  \label{tab:appendix_tfidf_lr_svm}
  \begin{tabular}{l l l}
    \toprule
    \textbf{Nhóm} & \textbf{Siêu tham số} & \textbf{Giá trị} \\
    \midrule
    \multirow{4}{*}{TF-IDF}
      & \texttt{max\_features}     & \tfidfmaxfeatures{} \\
      & \texttt{ngram\_range}      & $(1, 2)$ \\
      & \texttt{min\_df / max\_df} & 5 / 0{,}95 \\
      & \texttt{sublinear\_tf}     & True \\
    \midrule
    \multirow{4}{*}{LR}
      & \texttt{C} (tốt nhất)      & \lrbestc{} (solver=saga, L2) \\
      & \texttt{max\_iter}         & 3000 \\
      & \texttt{class\_weight}     & balanced \\
      & Lưới tìm kiếm             & $C \in \{0{,}5;\ 1;\ 5;\ 10;\ 20;\ 50\}$ \\
    \midrule
    \multirow{3}{*}{SVM}
      & \texttt{C} (tốt nhất)      & \svmbestc{} (LinearSVC + Platt) \\
      & \texttt{class\_weight}     & balanced \\
      & Lưới tìm kiếm             & $C \in \{0{,}1;\ 0{,}5;\ 1;\ 2;\ 5;\ 10\}$ \\
    \bottomrule
  \end{tabular}
\end{table}

\begin{table}[H]
  \centering
  \small
  \caption[Cấu hình BiLSTM và PhoBERT]{Cấu hình BiLSTM và PhoBERT.}
  \label{tab:appendix_dl}
  \begin{tabular}{l l l}
    \toprule
    \textbf{Mô hình} & \textbf{Siêu tham số} & \textbf{Giá trị} \\
    \midrule
    \multirow{7}{*}{BiLSTM}
      & Embedding dim            & \bilstmembdim{} (FastText cc.vi.300) \\
      & Hidden dim / Layers      & \bilstmhiddendim{} / \bilstmlayers{} \\
      & Max seq length           & \bilstmmaxseq{} \\
      & Dropout / Weight decay   & 0{,}3 / $10^{-4}$ \\
      & Learning rate            & $10^{-3}$ (AdamW) \\
      & Batch / Epochs           & 64 / 20 (patience=5) \\
      & Gradient clip            & $\max\_norm = 1{,}0$ \\
    \midrule
    \multirow{7}{*}{PhoBERT}
      & \texttt{model\_name}     & \texttt{vinai/phobert-base} \\
      & Max seq length           & \phobertmaxseq{} \\
      & Dropout                  & 0{,}1 \\
      & LR / Weight decay        & $3 \times 10^{-5}$ / 0{,}01 \\
      & Batch / Grad accum.      & 16 / 4 (effective $\approx 64$) \\
      & Epochs / Patience        & 8 / 2 \\
      & Layer LR decay           & 0{,}95 \\
    \bottomrule
  \end{tabular}
\end{table}

\subsection*{B.1. Knowledge Distillation}

Mô hình sinh viên là \texttt{StudentBiLSTM} --- một phiên bản BiLSTM rút gọn với \texttt{embedding\_dim = \studentembdim{}}, \texttt{hidden\_dim = \studenthiddim{}}, \texttt{num\_layers = \studentlayers{}}. Trọng số soft target $\alpha = \studentalpha{}$, nhiệt độ phân phối $T = \studenttemp{}$. Optimizer sử dụng AdamW với \texttt{lr}=$10^{-3}$, scheduler CosineAnnealingLR. Huấn luyện tối đa \studentepochs{} epoch với early stopping (\texttt{patience}=\studentpatience{}). Hàm mất mát KD chuẩn (Hinton et al., 2015) đã được trình bày ở Mục~\ref{sec:extended_analyses}; mô hình sinh viên được huấn luyện trên tập logits giáo viên đã được tính trước và lưu lại nhằm tiết kiệm chi phí tính toán.

\subsection*{B.2. Post-hoc Calibration}

Ba phương pháp hiệu chỉnh xác suất --- đều được khớp trên tập validation và áp dụng lên tập kiểm tra:

\begin{itemize}[leftmargin=2em,itemsep=2pt]
  \item \textbf{Platt scaling:} hồi quy logistic một chiều trên logits.
  \item \textbf{Temperature scaling:} tối ưu vô hướng $T \in [10^{-3},\ 10^{3}]$ bằng thuật toán L-BFGS-B.
  \item \textbf{Isotonic regression:} ánh xạ đơn điệu phi tham số trên $\sigma(z)$.
\end{itemize}

Giá trị $T$ đã khớp tương ứng: LR $=\lrtfit{}$, SVM $=\svmtfit{}$, BiLSTM $=\bilstmtfit{}$, PhoBERT $=\phoberttfit{}$.

\section*{Phụ lục C: Kết quả chi tiết}

\subsection*{C.1. Metric tổng hợp trên tập test}

\begin{table}[H]
  \centering
  \small
  \caption[Metric tổng hợp trên tập test]{Accuracy, Precision, Recall và F1-macro trên tập test.}
  \label{tab:appendix_full_metrics}
  \begin{tabular}{l c c c c}
    \toprule
    \textbf{Mô hình} & \textbf{Accuracy} & \textbf{Precision} & \textbf{Recall} & \textbf{F1-macro} \\
    \midrule
    Logistic Regression & \lracctab{} & \lrprectab{} & \lrrectab{} & \lrfonetab{} \\
    SVM                 & \svmacctab{} & \svmprectab{} & \svmrectab{} & \svmfonetab{} \\
    BiLSTM              & \bilstmacctab{} & \bilstmprectab{} & \bilstmrectab{} & \bilstmfonetab{} \\
    PhoBERT             & \phobertacctab{} & \phobertprectab{} & \phobertrectab{} & \phobertfonetab{} \\
    \bottomrule
  \end{tabular}
\end{table}

\subsection*{C.2. Metric theo lớp}

\begin{table}[H]
  \centering
  \small
  \caption[Precision/Recall/F1 theo lớp]{Precision, Recall và F1 theo lớp (Thật / Giả) trên tập test.}
  \label{tab:appendix_per_class}
  \begin{tabular}{l c c c c c c}
    \toprule
    & \multicolumn{3}{c}{\textbf{Lớp Thật (0)}} & \multicolumn{3}{c}{\textbf{Lớp Giả (1)}} \\
    \cmidrule(lr){2-4} \cmidrule(lr){5-7}
    \textbf{Mô hình} & Prec & Rec & F1 & Prec & Rec & F1 \\
    \midrule
    Logistic Regression & \lrrealprec{} & \lrrealrec{} & \lrrealfone{} & \lrfakeprec{} & \lrfakerec{} & \lrfakefone{} \\
    SVM                 & \svmrealprec{} & \svmrealrec{} & \svmrealfone{} & \svmfakeprec{} & \svmfakerec{} & \svmfakefone{} \\
    BiLSTM              & \bilstmrealprec{} & \bilstmrealrec{} & \bilstmrealfone{} & \bilstmfakeprec{} & \bilstmfakerec{} & \bilstmfakefone{} \\
    PhoBERT             & \phobertrealprec{} & \phobertrealrec{} & \phobertrealfone{} & \phobertfakeprec{} & \phobertfakerec{} & \phobertfakefone{} \\
    \bottomrule
  \end{tabular}
\end{table}

\subsection*{C.3. Confusion matrix trên tập test}

\begin{table}[H]
  \centering
  \small
  \caption[Confusion matrix trên tập test]{Số mẫu TN, FP, FN, TP của bốn mô hình trên tập test.}
  \label{tab:appendix_confusion}
  \begin{tabular}{l r r r r r}
    \toprule
    \textbf{Mô hình} & \textbf{TN} & \textbf{FP} & \textbf{FN} & \textbf{TP} & \textbf{Tổng lỗi} \\
    \midrule
    Logistic Regression & 985  & \lrfpcount{}   & \lrfncount{}   & 992 & \lrerrors{} (\lrerrorrate\%) \\
    SVM                 & 1.010 & \svmfpcount{}  & \svmfncount{}  & 756 & \svmerrors{} (\svmerrorrate\%) \\
    BiLSTM              & 977  & \bilstmfpcount{} & \bilstmfncount{} & 751 & \bilstmerrors{} (\bilstmerrorrate\%) \\
    PhoBERT             & \phoberttN{} & \phobertfP{} & \phobertfN{} & \phoberttP{} & \phoberterrors{} (\phoberterrorrate\%) \\
    \bottomrule
  \end{tabular}
\end{table}

\subsection*{C.4. Khoảng tin cậy Bootstrap 95\%}

Khoảng tin cậy được đo bằng phương pháp bootstrap với 10{,}000 lần lặp.

\begin{table}[H]
  \centering
  \small
  \caption[Khoảng tin cậy Bootstrap 95\% cho F1-score]{Khoảng tin cậy 95\% cho F1-score trên tập test, ước lượng bằng bootstrap với 10{,}000 lần lặp.}
  \label{tab:appendix_bootstrap}
  \begin{tabular}{l c}
    \toprule
    \textbf{Mô hình} & \textbf{95\% CI (F1-score)} \\
    \midrule
    Logistic Regression & \lrfonecI{} \\
    SVM                 & \svmfonecI{} \\
    BiLSTM              & \bilstmfonecI{} \\
    PhoBERT             & \phobertfonecI{} \\
    \bottomrule
  \end{tabular}
\end{table}

\subsection*{C.5. Knowledge Distillation -- chi tiết}

\begin{table}[H]
  \centering
  \small
  \caption[So sánh giáo viên -- sinh viên]{So sánh hai mô hình giáo viên (PhoBERT, BiLSTM) và sinh viên BiLSTM ($\alpha=\studentalpha$, $T=\studenttemp$).}
  \label{tab:appendix_distillation}
  \begin{tabular}{l r r r}
    \toprule
    \textbf{Mô hình} & \textbf{F1 test} & \textbf{Params} & \textbf{Kích thước (MB)} \\
    \midrule
    PhoBERT (giáo viên)  & \phobertfonetab{} & \phobertparamcount{} & \phobertsizemB{} \\
    BiLSTM (giáo viên)   & \bilstmfonetab{}  & \bilstmparamcount{}  & 24{,}57 \\
    Student BiLSTM (KD)  & \studentfone{}    & \studentparams{}    & \studentsizemB{} \\
    \bottomrule
  \end{tabular}
\end{table}

\noindent Tỷ lệ nén so với PhoBERT là \studentcompressionphi{$\times$}; so với BiLSTM là \studentcompressionbI{$\times$}. Tốc độ suy diễn đạt \studentlat{}~ms, nhanh hơn \studentspeedup{$\times$}; mô hình sinh viên giữ được \studentretainbI{} phần trăm chất lượng của BiLSTM.

\subsection*{C.6. Attribution Faithfulness}

Faithfulness là mức giảm xác suất trung bình khi che \texttt{top-5} token được attribution đánh dấu trên tập \nattrexamples{} mẫu.

\begin{table}[H]
  \centering
  \small
  \caption[Faithfulness của sáu phương pháp attribution]{Faithfulness của sáu phương pháp attribution.}
  \label{tab:appendix_faithfulness}
  \begin{tabular}{l c}
    \toprule
    \textbf{Phương pháp} & \textbf{Faithfulness} \\
    \midrule
    LR -- SHAP                     & \attrfaithfulnesslR{} \\
    SVM -- SHAP                    & \attrfaithfulnesssVM{} \\
    BiLSTM -- Integrated Gradients & \attrfaithfulnessbilstm{} \\
    BiLSTM -- Simple Gradients     & \attrfaithfulnessbilstmgrad{} \\
    PhoBERT -- SHAP                & \attrfaithfulnessphobERTSHAP{} \\
    PhoBERT -- Integrated Gradients& \attrfaithfulnessphobERTIG{} \\
    PhoBERT -- Attention Rollout   & \attrfaithfulnessphobERTRollout{} \\
    \bottomrule
  \end{tabular}
\end{table}

\subsection*{C.7. Cross-validation 5-fold}

\begin{table}[H]
  \centering
  \small
  \caption[Cross-validation 5-fold]{Accuracy và F1 trung bình qua 5-fold cross-validation (3 seed) cho hai mô hình truyền thống.}
  \label{tab:appendix_cv}
  \begin{tabular}{l c c c c}
    \toprule
    \textbf{Mô hình} & \textbf{Accuracy} & \textbf{$\pm$std} & \textbf{F1-macro} & \textbf{$\pm$std} \\
    \midrule
    Logistic Regression & \cvlRAcc{} & \cvlRAccstd{} & \cvlRFone{} & \cvlRFonestd{} \\
    SVM                 & \cvsVMAcc{} & \cvsVMAccstd{} & \cvsVMFone{} & \cvsVMFonestd{} \\
    \bottomrule
  \end{tabular}
\end{table}

\subsection*{C.8. Hard cases -- phân bố}

Tập \nhardexamples{} mẫu khó (cả bốn mô hình đều phân loại sai) được phân bố theo lý do thất bại suy đoán như sau:

\begin{table}[H]
  \centering
  \small
  \caption[Phân bố hard cases theo lý do]{Phân bố \nhardexamples{} mẫu khó theo lý do thất bại suy đoán.}
  \label{tab:appendix_hard_reasons}
  \begin{tabular}{l r}
    \toprule
    \textbf{Lý do} & \textbf{Số mẫu} \\
    \midrule
    \texttt{dataset\_ambiguity}      & \nhardambiguity{} \\
    \texttt{knowledge\_verification} & \nhardknowledge{} \\
    \texttt{semantic\_reasoning}     & \nhardsemantic{} \\
    \midrule
    \textbf{Chủ đề politics}        & \nhardpolitics{} (phân bố theo chủ đề) \\
    \textbf{Độ dài short ($<$30 từ)} & \nhardshort{} \\
    \bottomrule
  \end{tabular}
\end{table}

\section*{Phụ lục D: Tái lập thực nghiệm}

\subsection*{D.1. Yêu cầu môi trường}

\noindent\begin{tabular}{@{}l@{\;}l@{}}
\textbf{Python}     & \texttt{numpy 2.4}, \texttt{pandas 3.0}, \texttt{scikit-learn 1.8}, \\
                    & \texttt{scipy 1.17}, \texttt{torch 2.10}, \texttt{transformers 4.57}, \\
                    & \texttt{sentencepiece 0.2}, \texttt{fasttext-wheel 0.9.2}. \\
\textbf{Tiếng Việt} & \texttt{py\_vncorenlp 0.1.4} (Java $\geq$ 11) \\
                    & hoặc \texttt{underthesea 9.1} (dự phòng). \\
\textbf{Giải thích} & \texttt{shap $\geq$ 0.45}, \texttt{captum $\geq$ 0.7}. \\
\end{tabular}

Biến môi trường cần thiết: \texttt{JAVA\_HOME} trỏ tới JRE. PhoBERT (\texttt{vinai/phobert-base}) và FastText (\texttt{cc.vi.300.bin}) được tải tự động từ HuggingFace và lưu tại thư mục \texttt{data/fasttext/}.

\subsection*{D.2. Pipeline huấn luyện -- đánh giá}

\begin{verbatim}
python -m venv .venv && source .venv/bin/activate
pip install -r requirements.txt
python -m src.preprocessing.word_segmentation
python -m src.preprocessing.split_data
python -m src.features.extract_all_features
python -m src.training.train_all
python -m src.training.reproduce_predictions
python -m src.training.train_student
python -m src.evaluation.post_hoc_calibration
python -m src.evaluation.error_analysis
python -m src.evaluation.hard_cases
python -m src.analysis.statistical_tests
python -m src.analysis.generate_paper_tables
python -m src.analysis.generate_paper_figures
cd paper && ./build.sh
\end{verbatim}

\subsection*{D.3. Thiết lập ngẫu nhiên}

Seed chính của toàn bộ pipeline là \texttt{cfg.RANDOM\_STATE = \randomstate{}} --- được sử dụng cho \texttt{train\_test\_split}, \texttt{GridSearchCV} (qua \texttt{StratifiedKFold(shuffle=True)}) và các script huấn luyện. Hàm \texttt{set\_reproducibility\_seeds} cố định đồng thời các seed của \texttt{random}, \texttt{numpy}, \texttt{torch} và \texttt{torch.cuda}; hàm này được gọi ngay đầu mỗi trainer nhằm đảm bảo kết quả hoàn toàn tái lập.

\subsection*{D.4. File cấu hình}

Toàn bộ siêu tham số được tập trung tại file cấu hình chung và import qua singleton \texttt{cfg}, bao gồm các khối: \texttt{PATHS}, \texttt{DATA} (tỉ lệ chia tập, \texttt{min\_word\_count = \minwordcount{}}), \texttt{TFIDF}, \texttt{LR}, \texttt{SVM}, \texttt{BILSTM}, \texttt{PHOBERT}, \texttt{ANALYSIS} (\texttt{bootstrap\_iterations = 10{,}000}, \texttt{significance\_level = 0{,}05}).

Mọi giá trị hiển thị trong bài báo được định nghĩa tại \texttt{data\_commands.sty} dưới dạng \texttt{\textbackslash newcommand}, cho phép cập nhật đồng bộ từ một vị trí duy nhất --- tránh sai lệch giữa các phần của tài liệu.

